\paragraph{Statement A.5.6 (Proved smoothed expansion and conditional two-sided oscillation).}
Let
\[
 \mathcal A(Y)=\sum_{A_z\leq Y}
       \#\{\mathfrak a:N\mathfrak a=m_z,\ \theta(\mathfrak a)\in I\},
 \qquad e(Y)=\frac{\mathcal A(Y)-\lambda_0cY/3}{\sqrt Y}.
\]
Fix a smooth nonincreasing function $w:[0,\infty)\to[0,1]$,
equal to $1$ near zero and to $0$ for sufficiently large argument.
Its derivative is supported in a compact subinterval of $(0,\infty)$.
Then, as $Y\to\infty$,
\begin{align}\label{A:fixed-smoothing-AP}
 \int_0^\infty e(Yv)\sqrt v\,(-w'(v))\,dv
 &=\frac14\sum_{j\geq1}p'_j
 \left(r_j\widehat w(s_j)Q_{j,I}Y^{it_j}
 +\overline{r_j\widehat w(s_j)Q_{j,I}}Y^{-it_j}\right)
 \nonumber\\
 &\quad+O_w(Y^{-1/25}).
\end{align}
The series converges absolutely and uniformly in $\log Y$.
It is a real uniformly almost-periodic function with mean zero.
The following is consequently a sufficient condition for the
requested two-sided oscillation:
\begin{equation}\label{A:exact-nonvanishing-condition}
 \text{there is a }t>0\text{ such that }\quad
 \mathcal C(t):=\sum_{j:t_j=t}p'_jQ_{j,I}\ne0.
\end{equation}
If \eqref{A:exact-nonvanishing-condition} holds, then
$R(X)-C_0X=\Omega_+(X^{1/2})$ and
$R(X)-C_0X=\Omega_-(X^{1/2})$.
The sum in \eqref{A:exact-nonvanishing-condition} is over the
complete eigenspace and is independent of its real orthonormal
basis. Nonvanishing is not asserted in this statement.

\begin{proof}
We give the additional estimates needed to pass to the sharp arc.
They also explain why the $O(H)$ zero-residue bound used in
Statement A.5.5 cannot simply be retained in this argument.

\emph{1. Uniform angular discrepancy.}
The proof of Statement A.4.4 gives the stronger intermediate fact
\begin{equation}\label{A:all-interval-discrepancy}
 D(V):=\sup_{J\subset\mathbb R/\mathbb Z}
 \left|\sum_{A_z\leq V}\sum_{N\mathfrak a=m_z}
       \bigl(\mathbf1_J(\theta(\mathfrak a))-|J|\bigr)\right|
 \ll_\epsilon V^{2/3+\epsilon}.
\end{equation}
Indeed its phase-aligned bound for
$\sum_{\ell\leq H}|W_\ell|/\ell$ is independent of the interval;
the weighted Erd\H{o}s--Tur\'an step and the two radial
brackets therefore apply to every $J$ with the same constants.
One can use half-open intervals, or control a point mass by
the same positive Fej\'er majorant. This avoids assuming any
uniform Pell estimate for a varying interval endpoint.
For every periodic absolutely continuous function $f$ of mean zero,
Stieltjes integration by parts consequently gives
\begin{equation}\label{A:BV-discrepancy}
 \left|\sum_{A_z\leq V}\sum_{\mathfrak a}f(\theta(\mathfrak a))\right|
 \leq D(V)\operatorname{Var}(f).
\end{equation}
The assertion follows first for step functions, by expressing
them through interval indicators, and then by approximation.

\emph{2. Canonical angular polynomials have bounded variation.}
We use a specified choice of the Selberg polynomials.
The sawtooth approximation in
\cite[Lemma 2.2]{WangZhang2023}, explicitly attributed there to
Vaaler's Theorem 18, states the following. For an integer $H\geq1$,
put
\[
 W(t)=\pi t(1-|t|)\cot(\pi t)+|t|\quad(0<|t|<1),
 \qquad W(0)=1,
\]
extend $W$ by zero outside $[-1,1]$, and set
\[
 \psi_H(x)=-\sum_{1\leq|\ell|\leq H}
       \frac{W(\ell/(H+1))}{2\pi i\ell}e^{2\pi i\ell x},
 \quad
 F_H(x)=\sum_{|\ell|\leq H}
       \left(1-\frac{|\ell|}{H+1}\right)e^{2\pi i\ell x}.
\]
Then $F_H\geq0$ and
$|\psi_H-\psi|\leq F_H/(2(H+1))$.
These hypotheses are purely trigonometric and impose no
condition on our lattice. For an interval $I=[a,b]$ take
\[
 B_H^\pm(x)=b-a+\psi_H(a-x)+\psi_H(x-b)
 \ \pm\frac{F_H(a-x)+F_H(x-b)}{2(H+1)}.
\]
With the usual half-value convention at the two endpoints,
these satisfy the pointwise brackets of Statement A.5.5.
The endpoint convention changes our actual count by
$O(\log(2Y))$, already proved in Statement A.1.3.

In addition to the previously used Fourier bounds, these
particular polynomials satisfy
\begin{equation}\label{A:Selberg-BV}
 \|B_H^\pm\|_\infty+\operatorname{Var}(B_H^\pm)\ll1,
 \qquad \|(B_H^\pm)^{(r)}\|_1\ll_r H^{r-1}\quad(r\geq1).
\end{equation}
Here is a direct verification of the extra assertion.
The extension of $W$ is $C^1$, with integrable piecewise second
derivative: its value and first derivative vanish at $\pm1$.
Twice summing its Fourier polynomial by parts gives
\[
 \left|\sum_{\ell\in\mathbb Z}W(\ell/(H+1))e^{2\pi i\ell x}\right|
 \ll \min\{H,(H\|x\|^2)^{-1}\}.
\]
Its $L^1$ norm is therefore bounded independently of $H$.
Since $\psi_H'$ equals one minus this polynomial, it has
bounded $L^1$ norm. The formula
$F_H(x)=\sin^2(\pi(H+1)x)/((H+1)\sin^2(\pi x))$
shows $\|F_H'\|_1\ll H$: split at $\|x\|=1/H$ and
differentiate, using $O(H^3|x|)$ in the first range and
$O(|x|^{-2}+H^{-1}|x|^{-3})$ in the second.
This proves the variation bound; bounded sup norm follows from
the sawtooth approximation and $F_H\leq H+1$.
Finally, the elementary $L^1$ Bernstein bound for band-limited
polynomials gives the derivative estimates in
\eqref{A:Selberg-BV}. It can be proved here without an
approximation theorem: choose a fixed smooth frequency cutoff
equal to one on $[-1,1]$ and supported in $[-2,2]$.
Its periodic kernel, differentiated $r-1$ times, has $L^1$
norm $O_r(H^{r-1})$, by periodizing its Schwartz inverse Fourier
transform. Convolve this kernel with $(B_H^\pm)'$.

\emph{3. Improvement of the zero residue.}
For either polynomial $B=B_H^\pm$, the exact identity at $s=0$
in Statement A.5.5 gives
\[
 H_B(0):=\sum_{|\ell|\leq H}b_\ell H_\ell(0)
   =\sum_z\sum_{\mathfrak a}f_{A_z}(\theta(\mathfrak a)),
\]
where, putting $\delta(A)=\operatorname{arsinh}(1/A)/L_0$,
\[
 f_A(x)=\tfrac12\{B(x+\delta(A))+B(x-\delta(A))-2B(x)\}.
\]
This series converges absolutely for each fixed $H$,
since its large-$A$ summand is $O_H(A^{-2})$ and
$S_A(V)\ll V$. Every $f_A$ has mean zero.
On a dyadic block $A\asymp Z$, \eqref{A:Selberg-BV} gives
\[
 \operatorname{Var}(f_A)\ll
      \min\{1,H^2/Z^2\},\qquad
 \operatorname{Var}(\partial_A f_A)
      \ll\min\{H/Z^2,H^2/Z^3\}.
\]
The second estimate follows from
$|\delta'(A)|\ll A^{-2}$ and the first or second difference
of $B''$; the first follows from the second difference of $B'$.
Partial summation in $A$, followed by
\eqref{A:BV-discrepancy}, bounds that dyadic block by
\[
 \begin{cases}
 O_\epsilon(HZ^{\vartheta-1}),&Z\leq H,\\
 O_\epsilon(H^2Z^{\vartheta-2}),&Z\geq H,
 \end{cases}
 \qquad \vartheta=2/3+\epsilon<1.
\]
For $A\leq V$ use instead $\|B\|_\infty\ll1$ and
$S_A(V)\ll V$. Summing the remaining dyadic blocks gives
\[
 |H_B(0)|\ll_\epsilon
 V+HV^{\vartheta-1}+H^\vartheta.
\]
Take $V=H^{1/(2-\vartheta)}$. Since $\epsilon$ can be
arbitrarily small, this proves
\begin{equation}\label{A:zero-residue-improved}
 H_B(0)\ll_\epsilon H^{3/4+\epsilon}.
\end{equation}

\emph{4. Fixed radial smoothing and passage to the sharp arc.}
Use the negative-line transform of Statement A.5.5 with the fixed
function $w$, $\sigma=-49/100$, and $H=\lfloor Y^{3/5}\rfloor$.
The parameter $T$ in the transform bounds is now a fixed
constant depending on $w$. The contour bound remains
$O_w((T+H)\log(2H))$ for $-1/2<\sigma<0$.
To verify the only extra range not used in Statement A.5.5,
take $K>T$ in its high-angular calculation. The dyadic sum is
\[
 K^{1/2-\sigma}\sum_{J\ll K}J^\sigma\min(J,T)^{1/2}
 \ll_\sigma K^{1/2-\sigma}T^{\sigma+1/2}\leq C_\sigma K.
\]
The low and comparable ranges use only $a<1$, exactly as before.
The spherical-correction inversion, the new zero-residue bound,
and the exceptional intersection terms therefore leave an error
\[
 O_w\bigl(Y^\sigma H\log(2H)
          +Y^\sigma H^{1-\sigma}\log(2H)
          +H^{3/4+1/100}+\log(2Y)\bigr)
   =O_w(Y^{23/50}).
\]
The zero Fourier mode has the same main term and pole at $s=1$
as in Statement A.5.5; all $M_\ell(0)$ vanish.

We also need a bound uniform in $H$ for the individual
$Q_{j,H}^\pm$. The quantitative assertion of
\cite[Theorem 7.1 and Section 7.1]{AnantharamanZelditch2012}
gives a polynomial bound in $t_j$ for the H\"older norm of
the boundary primitive. Its displayed constant is
$O((1+t_j)^{3/2})$ times
$\|u_j\|_\infty+\|\nabla u_j\|_\infty$.
On the fixed compact cover, the elliptic estimate and Sobolev
embedding give the deliberately coarse bound
$\|u_j\|_{C^1}\ll(1+t_j)^4$.
The fixed coordinate changes and multiplication by
$(2e^v)^{-s_j}$ in Statement A.4.6 cost at most a further fixed
power of $1+t_j$. Thus, for some absolute integer $M$ depending
only on this fixed geometric setup, the periodic primitive
$F_j$ of $\beta_j$ has
\[
 \|F_j\|_{C^{1/2}}\ll(1+t_j)^M.
\]
An affine mean is understood in this assertion, as in
Statement A.4.6.

Identify $u=L_0\theta$, absorbing only fixed coordinate factors.
The middle polynomial in $B_H^\pm$ is the convolution of the
interval indicator with the kernel having coefficients
$W(\ell/(H+1))$. Its bounded $L^1$ norm and the displayed
decay bound imply approximation of a $1/2$-H\"older primitive
with error $O(H^{-1/2})$. For the two Fej\'er terms,
integration by parts and the preceding bounds for $F_H'$ give
\[
 \frac1{H+1}\left|\int F_H(x-a)\,\beta_j(dx)\right|
 \ll (1+t_j)^M H^{-1/2};
\]
subtract $F_j(a)$ inside the derivative integral and use
$|F_j(x)-F_j(a)|\ll(1+t_j)^M\|x-a\|^{1/2}$.
Consequently
\begin{equation}\label{A:Selberg-Q-limit}
 |Q_{j,H}^\pm-Q_{j,I}|
       \ll(1+t_j)^M H^{-1/2},
 \qquad |Q_{j,H}^\pm|+|Q_{j,I}|\ll(1+t_j)^M,
\end{equation}
after increasing $M$ if necessary.
The rapid vertical decay of $\widehat w$, the polynomial gamma
bound on $r_j$, the broad-band period bound, and
$\#\{t_j\leq V\}\ll V^2$ make the residue series absolutely
convergent. Equation \eqref{A:Selberg-Q-limit} permits passage
to its sharp-arc limit with error $O_w(Y^{1/2}H^{-1/2})$.

Bracket the actual radially smoothed arc count by $B_H^-$ and
$B_H^+$. Their main terms differ from
$\lambda_0cY\widehat w(1)/3$ by $O_w(Y/H)$.
Both spectral terms converge to the same series, and all the
other errors just proved are $O_w(Y^{23/50})$.
This proves the corresponding unnormalized expansion.
Partial summation gives
\[
 \sum_z\sum_{\mathfrak a}w(A_z/Y)\mathbf1_I(\theta(\mathfrak a))
       -\frac{\lambda_0c}{3}Y\widehat w(1)
 =\sqrt Y\int_0^\infty e(Yv)\sqrt v\,(-w'(v))\,dv,
\]
and proves \eqref{A:fixed-smoothing-AP}.

\emph{5. The conditional two-sided conclusion.}
Group the absolutely convergent series by distinct $t$.
Its coefficient at frequency $t$ is
$r(t)\widehat w(1/2+it)\mathcal C(t)/4$.
Here $r(t)\ne0$ because the gamma function has no zeros.
If \eqref{A:exact-nonvanishing-condition} holds, choose $w$
with $\widehat w(1/2+it)\ne0$: smooth monotone approximations
to $\mathbf1_{[0,1]}$ have Mellin transforms tending to
$(1/2+it)^{-1}\ne0$.
The resulting real series $P(\log Y)$ is nonzero.
Termwise averaging, justified by absolute uniform convergence,
gives mean zero and strictly positive mean square.
Since it is bounded, these two facts imply
$\limsup P>0$ and $\liminf P<0$.
For example, if $\limsup P\leq0$, then for every $\varepsilon>0$
one has eventually $P\leq\varepsilon$ and $P\geq-C$.
The inequality
$P^2\leq(\varepsilon-C)P+\varepsilon C$ contradicts the
positive mean square after averaging and letting
$\varepsilon\to0$. The other sign is identical.

The kernel $\sqrt v\,(-w'(v))$ is nonnegative, has finite
positive integral, and is supported in a fixed compact
subinterval of $(0,\infty)$. Positive and negative values
bounded away from zero of its convolution with $e(Y)$
therefore force positive and negative values bounded away
from zero of $e$ at arbitrarily large arguments.
Finally $Y=\sqrt{9X^2-1}=3X+O(X^{-1})$ and
$R(X)=\mathcal A(Y)+O(\log(2X))$ transfer both assertions
to $R(X)-C_0X$. The eigenspace sum is basis-independent because
it is the contraction of the two linear period functionals
with the orthogonal projector onto that eigenspace.
\end{proof}
